%% ma: L12-B08-0171
%% lop: 12
%% chuong: 2
%% bai: 8
%% dang: 12.8.53
%% loai: TN4
%% muc: VD
%% dap_an: "B"
%% nguon: "CĐ2.4.pdf giáo viên gửi 27/10/2026 (M07, Lê Minh Khôi – Nguyễn Anh Khôi, Chương II Chủ đề 4 Các bài toán về tâm tỉ cự), Câu 21"
%% kien_thuc: "Bài 8 (khoảng cách hai điểm, biểu thức tọa độ), Bài 6 (bình phương vô hướng, tâm tỉ cự)"
%% trang_thai: chinh-thuc
%% ghi_chu: "Giáo viên đã duyệt 27/10/2026. Đề gốc dùng chung giả thiết cho Câu 20–22; đã chép vào từng câu. Phương án D gốc ghi nhầm ký hiệu K_min, đã sửa thành B_min"
\begin{ex}
Trong không gian $Oxyz$, cho ba điểm $M(2;1;-2)$, $N(-1;0;3)$, $P(1;-3;5)$ và điểm $I$ bất kỳ trong không gian. Tìm giá trị nhỏ nhất của biểu thức $B=2IM^2-IN^2+5IP^2$.
\choice{$B_{\min}=\dfrac{505}2$}{\True $B_{\min}=\dfrac{505}6$}{$B_{\min}=\dfrac{101}2$}{$B_{\min}=\dfrac{55}6$}
\loigiai{Gọi $E$ thỏa $2\overrightarrow{EM}-\overrightarrow{EN}+5\overrightarrow{EP}=\vec 0\Rightarrow E=\dfrac{2M-N+5P}6=\left(\dfrac53;-\dfrac{13}6;3\right)$.

$B=6IE^2+2EM^2-EN^2+5EP^2\ge\dfrac{505}6$.}
\end{ex}
